\chapter{Abstract}
\begin{sloppy}
This thesis presents a study on the ferromagnetic resonance (FMR) of symmetric heavy metal-\allowbreak superconductor-\allowbreak ferromagnet-\allowbreak superconductor-\allowbreak heavy metal (HM-SC-FM-SC-HM) stacks, specifically focusing on $\mathrm{Pt}/\allowbreak\mathrm{Nb}/\allowbreak{\mathrm{Ni}}_{80}{\mathrm{Fe}}_{20}/\allowbreak\mathrm{Nb}/\allowbreak \mathrm{Pt}$ heterostructures. The primary objective is to investigate the spin pumping through the superconductor, particularly near the critical temperature $T_\textrm{C}$ of Niobium (Nb).
\end{sloppy}

Using a coplanar waveguide (CPW) setup, FMR measurements were conducted on samples with varying Nb thicknesses $d_\textrm{Nb}$ of 45 nm and 30 nm. The study examines key parameters such as the Gilbert damping constant $\alpha$, the inhomogeneous broadening $\Delta H_0$, the effective magnetization $M_\textrm{eff}$, and the resonance field shift $H_\textrm{shift}$ as functions of temperature. The results indicate a significant spiking behavior in all parameters near $T_\textrm{C}$, suggesting a disturbance in the spin-transfer mechanism, potentially due to domain formation or vortex motion in the superconductor. The findings for the $d_\textrm{Nb} = \SI{45}{nm}$ sample align well with existing literature, while the $d_\textrm{Nb} = \SI{30}{nm}$ sample exhibits deviations that may be attributed to sample aging.

Additionally, the thesis details the fabrication of a device capable of FMR measurements under a DC-bias current using on-chip waveguides. The microfabrication process involved the use of optical lithography and the development of an isolation layer between the CPW and metallic film. For this, a process protocol was established, and a test sample was fabricated to validate the approach.